% !TeX program = pdflatex
% Tài liệu độc lập, tương thích với trình biên dịch pdfLaTeX mặc định của Overleaf.
\documentclass[12pt,a4paper]{article}

\usepackage[a4paper,top=2.2cm,bottom=2.2cm,left=2.5cm,right=2cm]{geometry}
\usepackage[utf8]{inputenc}
\usepackage[T5]{fontenc}
\usepackage[vietnamese]{babel}
\usepackage{lmodern}

\usepackage{microtype}
\usepackage{setspace}
\usepackage{enumitem}
\usepackage{array}
\usepackage{booktabs}
\usepackage{longtable}
\usepackage{tabularx}
\usepackage{multirow}
\usepackage{makecell}
\usepackage{ragged2e}
\usepackage{graphicx}
\usepackage{caption}
\usepackage{float}
\usepackage{pdflscape}
\usepackage[table]{xcolor}
\usepackage{fancyhdr}
\usepackage{hyperref}

\hypersetup{
  colorlinks=true,
  linkcolor=black,
  urlcolor=blue,
  pdftitle={Kịch bản thử nghiệm hệ thống quản lý SBOM},
  pdfauthor={Nhóm phát triển}
}

\setstretch{1.15}
\setlength{\parindent}{1.2cm}
\setlength{\parskip}{3pt}
\setlength{\tabcolsep}{4pt}
\renewcommand{\arraystretch}{1.18}
\setlist[itemize]{leftmargin=1.7cm,itemsep=2pt,topsep=3pt}
\setlist[enumerate]{leftmargin=1.7cm,itemsep=2pt,topsep=3pt}

\definecolor{headergray}{RGB}{230,230,230}
\definecolor{lightgreen}{RGB}{226,239,218}
\definecolor{lightred}{RGB}{255,230,230}

\pagestyle{fancy}
\fancyhf{}
\fancyhead[L]{\small Kịch bản thử nghiệm hệ thống quản lý SBOM}
\fancyhead[R]{\small \thepage}
\renewcommand{\headrulewidth}{0.4pt}

\captionsetup{font=small,labelfont=bf,labelsep=period}
\renewcommand{\figurename}{Hình}
\renewcommand{\tablename}{Bảng}

\newcolumntype{Y}{>{\RaggedRight\arraybackslash}X}
\newcolumntype{C}[1]{>{\Centering\arraybackslash}p{#1}}
\newcolumntype{L}[1]{>{\RaggedRight\arraybackslash}p{#1}}
\newcommand{\code}[1]{\texttt{#1}}
\newcommand{\PASS}{\cellcolor{lightgreen}\textbf{Đạt}}
\newcommand{\FAIL}{\cellcolor{lightred}\textbf{Không đạt}}

% Hiển thị ảnh nếu tồn tại; nếu chưa có thì tạo khung giữ chỗ.
\newcommand{\evidencefigure}[3][0.88\textwidth]{%
  \begin{figure}[H]
    \centering
    \IfFileExists{#2}{%
      \includegraphics[width=#1,height=0.58\textheight,keepaspectratio]{#2}%
    }{%
      \fbox{\parbox[c][5.2cm][c]{#1}{%
        \centering\small Chèn ảnh minh chứng tại đây\\[3pt]
        \code{#2}}}%
    }
    \caption{#3}
  \end{figure}
}

\begin{document}

\section{Mục tiêu thử nghiệm}

Tài liệu này mô tả các kịch bản thử nghiệm đối với nguyên mẫu hệ thống quản lý
Software Bill of Materials (SBOM). Hoạt động thử nghiệm nhằm xác nhận các mục tiêu sau:

\begin{itemize}
  \item Hệ thống có thể tiếp nhận và phân tích SBOM định dạng CycloneDX JSON.
  \item Hệ thống có thể phân tích repository, phát hiện tệp khai báo phụ thuộc và sinh SBOM bằng Syft.
  \item Hệ thống có thể đối chiếu SBOM đầu vào với inventory sinh lại từ mã nguồn.
  \item Các sai lệch được phân loại đúng thành \code{MATCHED},
        \code{MISSING\_IN\_SBOM}, \code{EXTRA\_IN\_SBOM} và
        \code{VERSION\_MISMATCH}.
  \item Hệ thống lưu được validation run, snapshot, pipeline run và bằng chứng liên quan.
  \item Grype có thể làm giàu kết quả bằng thông tin lỗ hổng và hệ thống phân biệt được
        quét thành công với quét thất bại.
  \item Kết quả kiểm chứng có thể được xuất thành workbook Excel phục vụ lưu trữ bằng chứng.
\end{itemize}

\textbf{Nguyên tắc đánh giá:} kết quả của luồng chức năng và kết quả nội dung SBOM
được đánh giá riêng. Một kịch bản chức năng được xem là \emph{đạt} khi hệ thống hoàn
thành đúng quy trình và trả về kết quả mong đợi. Nội dung SBOM chỉ \code{PASS} khi
không tồn tại component thiếu, thừa hoặc sai phiên bản. Vì vậy, luồng chức năng có
thể đạt trong khi nội dung SBOM là \code{FAIL}; trường hợp này chứng tỏ chức năng
phát hiện sai lệch hoạt động đúng.

\section{Phạm vi và đối tượng thử nghiệm}

\subsection{Các thành phần thuộc phạm vi}

\begin{itemize}
  \item Giao diện web React/TypeScript dùng để quản lý dự án, tải lên hoặc sinh SBOM,
        xem kết quả và xuất báo cáo.
  \item Backend Node.js/Express/TypeScript cung cấp REST API và xử lý nghiệp vụ.
  \item PostgreSQL lưu dự án, SBOM, component, dependency, vulnerability, snapshot,
        pipeline run và validation run.
  \item Syft 1.45.1 sinh CycloneDX JSON từ repository hoặc thư mục mã nguồn.
  \item Grype 0.114.0 quét SBOM để bổ sung thông tin CVE.
  \item ExcelJS 4.4.0 tạo báo cáo kiểm chứng định dạng \code{.xlsx}.
\end{itemize}

\subsection{Ngoài phạm vi}

Tài liệu không xác nhận mức sẵn sàng production, khả năng mở rộng ngang, hiệu năng tải
lớn, an toàn khai thác CVE trong môi trường thực tế, xác thực/phân quyền hoặc mức bao
phủ end-to-end hoàn chỉnh.

\section{Môi trường và điều kiện thử nghiệm}

\begin{table}[H]
  \centering
  \caption{Cấu hình môi trường thử nghiệm}
  \begin{tabularx}{\textwidth}{|L{3.7cm}|Y|}
    \hline
    \rowcolor{headergray}\textbf{Thành phần} & \textbf{Cấu hình} \\ \hline
    Frontend & Ứng dụng React/Vite, giao tiếp với backend qua REST API. \\ \hline
    Backend & Server Node.js/Express xử lý nghiệp vụ và gọi công cụ bên ngoài. \\ \hline
    Cơ sở dữ liệu & PostgreSQL lưu dữ liệu thử nghiệm và lịch sử chạy. \\ \hline
    Công cụ sinh SBOM & Syft 1.45.1 trên cùng môi trường với backend. \\ \hline
    Công cụ quét lỗ hổng & Grype 0.114.0; trạng thái scanner và lỗi được lưu riêng. \\ \hline
    Công cụ xuất báo cáo & ExcelJS 4.4.0. \\ \hline
    Phạm vi triển khai & Môi trường phát triển; pipeline là backend runner nội bộ. \\ \hline
  \end{tabularx}
\end{table}

\subsection{Điều kiện tiên quyết chung}

\begin{itemize}
  \item Frontend, backend và PostgreSQL đang hoạt động; frontend gọi được REST API.
  \item Syft và Grype đã được cài đặt, có quyền thực thi từ backend.
  \item Repository thử nghiệm có thể được tải về hoặc đã tồn tại tại đường dẫn hợp lệ.
  \item Người thử nghiệm có tệp CycloneDX JSON hợp lệ và các tệp SBOM đã được chỉnh sửa
        để tạo trường hợp thiếu, thừa hoặc sai phiên bản.
  \item Thư mục tạm của backend có đủ dung lượng; kết nối mạng khả dụng khi cần tải
        repository hoặc cơ sở dữ liệu lỗ hổng.
\end{itemize}

\section{Dữ liệu thử nghiệm}

Tại thời điểm chốt số liệu ngày 15/07/2026, cơ sở dữ liệu thử nghiệm ghi nhận
24 hệ thống, 30 pipeline, 16 pipeline run, 37 SBOM và 13 snapshot. Mô-đun kiểm
chứng ghi nhận 31 validation run trên 12 dự án; trong đó 14 run có báo cáo đối
chiếu SBOM với mã nguồn. Bốn run đại diện được chọn để đánh giá như sau.

\begin{table}[H]
  \centering
  \caption{Bộ dữ liệu đại diện}
  \begin{tabularx}{\textwidth}{|C{1.6cm}|L{3.5cm}|C{3.4cm}|Y|}
    \hline
    \rowcolor{headergray}
    \textbf{Mã} & \textbf{Dự án} & \textbf{Quy mô source scan} & \textbf{Mục đích} \\ \hline
    DATA01 & BookStack & 157 component, 110 dependency & SBOM gần khớp hoàn toàn với source. \\ \hline
    DATA02 & ThingsBoard & 2.281 component, 5.185 dependency & Kiểm tra repository quy mô lớn. \\ \hline
    DATA03 & Taxi Booking Web App & 355 component, 798 dependency & Minh họa nhiều nhóm mismatch và Trust Score trung bình. \\ \hline
    DATA04 & OFBiz Framework & 159 component, 140 dependency & Kiểm tra sai lệch lớn và Grype enrichment. \\ \hline
  \end{tabularx}
\end{table}

Trust Score được tính theo công thức:
\[
  \mathrm{TrustScore}=
  \frac{\mathrm{MatchedExact}}
       {\mathrm{SourceComponentCount}+\mathrm{ExtraInSbomCount}}
  \times 100.
\]

Trust Score phản ánh mức khớp inventory giữa SBOM đầu vào và SBOM sinh lại từ source;
đây không phải điểm an toàn bảo mật và không chứng minh một CVE có thể bị khai thác.

\section{Danh sách kịch bản thử nghiệm}

\begin{table}[H]
  \centering
  \caption{Tổng hợp các kịch bản thử nghiệm}
  \begin{tabularx}{\textwidth}{|C{2.1cm}|Y|L{3.2cm}|C{1.8cm}|}
    \hline
    \rowcolor{headergray}
    \textbf{Mã} & \textbf{Mục tiêu} & \textbf{Dữ liệu} & \textbf{Ưu tiên} \\ \hline
    SBOM-KT01 & Tải lên và phân tích CycloneDX JSON. & SBOM hợp lệ và không hợp lệ & Cao \\ \hline
    SBOM-KT02 & Phân tích repository, phát hiện manifest và sinh SBOM bằng Syft. & DATA01--DATA04 & Cao \\ \hline
    SBOM-KT03 & Đối chiếu SBOM với source và phân loại mismatch. & DATA01--DATA04 & Cao \\ \hline
    SBOM-KT04 & Tạo pipeline run, cập nhật step và lưu snapshot/kết quả. & Dự án có repository & Trung bình \\ \hline
    SBOM-KT05 & Xuất workbook Excel chứa tổng quan, diff và evidence. & Validation run hoàn tất & Cao \\ \hline
    SBOM-KT06 & Chạy Grype và phân biệt quét thành công với quét thất bại. & DATA04 và cấu hình lỗi & Cao \\ \hline
  \end{tabularx}
\end{table}

\section{Chi tiết kịch bản thử nghiệm}

\subsection{SBOM-KT01 -- Tải lên và phân tích CycloneDX JSON}

\begin{longtable}{|L{3.4cm}|L{11.2cm}|}
  \hline
  \rowcolor{headergray}\textbf{Thuộc tính} & \textbf{Nội dung} \\ \hline
  Mục tiêu & Xác nhận hệ thống tiếp nhận, kiểm tra và phân tích tệp CycloneDX JSON. \\ \hline
  Tiền điều kiện & Hệ thống hoạt động; người dùng đã chọn một dự án hợp lệ. \\ \hline
  Dữ liệu vào & Một tệp CycloneDX JSON hợp lệ; một tệp JSON sai cấu trúc để kiểm tra nhánh lỗi. \\ \hline
  Các bước thực hiện &
  1. Mở màn hình dự án và chọn chức năng tải SBOM.\newline
  2. Chọn tệp CycloneDX JSON hợp lệ và xác nhận tải lên.\newline
  3. Mở trang chi tiết SBOM, kiểm tra metadata, component và dependency.\newline
  4. Lặp lại với tệp JSON sai cấu trúc. \\ \hline
  Kết quả mong đợi &
  Tệp hợp lệ được lưu và dữ liệu được hiển thị đúng; tệp không hợp lệ bị từ chối với
  thông báo lỗi có thể hiểu được và không tạo bản ghi SBOM hỏng. \\ \hline
  Tiêu chí đạt & Tất cả dữ liệu chính được trích xuất; nhánh lỗi không làm gián đoạn dịch vụ. \\ \hline
\end{longtable}

\subsection{SBOM-KT02 -- Phân tích repository và sinh SBOM}

\begin{longtable}{|L{3.4cm}|L{11.2cm}|}
  \hline
  \rowcolor{headergray}\textbf{Thuộc tính} & \textbf{Nội dung} \\ \hline
  Mục tiêu & Xác nhận backend tải source, phát hiện manifest và gọi Syft để sinh CycloneDX JSON. \\ \hline
  Tiền điều kiện & Repository truy cập được; Syft 1.45.1 có quyền thực thi. \\ \hline
  Dữ liệu vào & URL repository hoặc đường dẫn source của DATA01--DATA04. \\ \hline
  Các bước thực hiện &
  1. Gắn repository với dự án.\newline
  2. Chọn chức năng phân tích source.\newline
  3. Kiểm tra danh sách tệp dependency manifest được phát hiện.\newline
  4. Chọn sinh SBOM và chờ quá trình hoàn thành.\newline
  5. Mở SBOM vừa tạo, kiểm tra metadata, component và dependency. \\ \hline
  Kết quả mong đợi & Repository được xử lý; Syft trả về CycloneDX JSON hợp lệ; dữ liệu
  được chuẩn hóa và lưu trong cơ sở dữ liệu. \\ \hline
  Tiêu chí đạt & Luồng hoàn tất không lỗi và SBOM sinh ra có số lượng component lớn hơn 0. \\ \hline
\end{longtable}

\subsection{SBOM-KT03 -- Đối chiếu SBOM với mã nguồn}

\begin{longtable}{|L{3.4cm}|L{11.2cm}|}
  \hline
  \rowcolor{headergray}\textbf{Thuộc tính} & \textbf{Nội dung} \\ \hline
  Mục tiêu & Xác nhận hệ thống phát hiện và phân loại chính xác sai lệch component. \\ \hline
  Tiền điều kiện & Source đã được phân tích; có SBOM cần kiểm chứng. \\ \hline
  Dữ liệu vào & DATA01--DATA04; các biến thể SBOM có component bị thêm, xóa hoặc đổi phiên bản. \\ \hline
  Các bước thực hiện &
  1. Mở chức năng kiểm chứng SBOM.\newline
  2. Chọn repository/source path và tải SBOM cần đối chiếu.\newline
  3. Chạy kiểm chứng, chờ hệ thống sinh inventory từ source.\newline
  4. Kiểm tra bốn nhóm kết quả và Trust Score.\newline
  5. Mở lại validation run để xác nhận kết quả đã được lưu. \\ \hline
  Kết quả mong đợi & Mỗi component được xếp đúng nhóm; số liệu tổng hợp khớp chi tiết;
  Trust Score đúng công thức; validation run lưu được report, graph và lỗi nếu có. \\ \hline
  Tiêu chí đạt & Chức năng hoàn thành và phát hiện đúng các sai lệch đã chủ động tạo. \\ \hline
\end{longtable}

\subsection{SBOM-KT04 -- Pipeline, snapshot và lịch sử chạy}

\begin{longtable}{|L{3.4cm}|L{11.2cm}|}
  \hline
  \rowcolor{headergray}\textbf{Thuộc tính} & \textbf{Nội dung} \\ \hline
  Mục tiêu & Xác nhận pipeline nội bộ tạo run, cập nhật step và lưu snapshot/kết quả. \\ \hline
  Tiền điều kiện & Dự án có repository hợp lệ và cấu hình pipeline. \\ \hline
  Dữ liệu vào & Một dự án có thể clone và sinh SBOM. \\ \hline
  Các bước thực hiện &
  1. Mở màn hình pipeline và kích hoạt một lần chạy.\newline
  2. Theo dõi trạng thái từng step.\newline
  3. Chờ run kết thúc và kiểm tra snapshot/change log.\newline
  4. Mở lại lịch sử pipeline sau khi tải lại trang. \\ \hline
  Kết quả mong đợi & Trạng thái chuyển đổi hợp lệ; thời điểm bắt đầu/kết thúc và lỗi được
  lưu; snapshot có component/dependency tương ứng và lịch sử không bị mất. \\ \hline
  Tiêu chí đạt & Pipeline run và các step có trạng thái cuối; kết quả truy xuất lại được. \\ \hline
\end{longtable}

\subsection{SBOM-KT05 -- Xuất báo cáo Excel}

\begin{longtable}{|L{3.4cm}|L{11.2cm}|}
  \hline
  \rowcolor{headergray}\textbf{Thuộc tính} & \textbf{Nội dung} \\ \hline
  Mục tiêu & Xác nhận hệ thống xuất đầy đủ kết quả và bằng chứng của validation run. \\ \hline
  Tiền điều kiện & Một validation run đã hoàn tất. \\ \hline
  Dữ liệu vào & Validation run đại diện từ DATA01--DATA04. \\ \hline
  Các bước thực hiện &
  1. Mở chi tiết validation run.\newline
  2. Chọn chức năng xuất báo cáo Excel.\newline
  3. Mở workbook vừa tải xuống.\newline
  4. Kiểm tra tên sheet, số liệu tổng quan, source diff, CVE và evidence. \\ \hline
  Kết quả mong đợi & Workbook mở được và có các sheet:
  \code{TongQuan}, \code{ThongKeCVE}, \code{CVEChiTiet},
  \code{KiemChungSource}, \code{PhanTichSource}, \code{TepPhuThuoc},
  \code{BangChung}, \code{QuyTrinhKiemThu}, \code{ChuGiaiCVE}. \\ \hline
  Tiêu chí đạt & Số liệu trong workbook khớp giao diện và validation run; không có sheet hỏng. \\ \hline
\end{longtable}

\subsection{SBOM-KT06 -- Grype enrichment và xử lý lỗi scanner}

\begin{longtable}{|L{3.4cm}|L{11.2cm}|}
  \hline
  \rowcolor{headergray}\textbf{Thuộc tính} & \textbf{Nội dung} \\ \hline
  Mục tiêu & Xác nhận Grype bổ sung finding và hệ thống không đồng nhất lỗi quét với 0 lỗ hổng. \\ \hline
  Tiền điều kiện & Có SBOM hợp lệ; Grype hoạt động cho nhánh thành công và được cấu hình
  sai có chủ đích cho nhánh thất bại. \\ \hline
  Dữ liệu vào & DATA04; một cấu hình đường dẫn Grype không hợp lệ hoặc lỗi có kiểm soát. \\ \hline
  Các bước thực hiện &
  1. Chạy Grype enrichment trên SBOM của DATA04.\newline
  2. Kiểm tra trạng thái scanner, CVE, severity, package và fixed version.\newline
  3. Lặp lại với cấu hình làm scanner thất bại.\newline
  4. Xuất workbook của cả hai trường hợp. \\ \hline
  Kết quả mong đợi & Nhánh thành công lưu finding; nhánh lỗi lưu trạng thái \code{FAILED}
  và thông báo lỗi. Báo cáo không diễn giải scanner thất bại thành ``không có lỗ hổng''. \\ \hline
  Tiêu chí đạt & Trạng thái, finding và lỗi được lưu/hiển thị đúng ở giao diện và Excel. \\ \hline
\end{longtable}

\section{Kết quả thử nghiệm}

\begin{table}[H]
  \centering
  \caption{Kết quả đối chiếu trên các dự án đại diện}
  \small
  \begin{tabularx}{\textwidth}{|L{3.1cm}|r|r|r|r|C{1.8cm}|L{1.7cm}|C{1.7cm}|}
    \hline
    \rowcolor{headergray}
    \textbf{Dự án} & \textbf{Matched} & \textbf{Missing} & \textbf{Extra} &
    \textbf{Version} & \textbf{Trust Score} & \textbf{CVE} & \textbf{Nội dung} \\ \hline
    BookStack & 157 & 1 & 1 & 0 & 98,74\% & Chưa quét & \FAIL \\ \hline
    ThingsBoard & 1.791 & 38 & 101 & 40 & 90,91\% & Chưa quét & \FAIL \\ \hline
    Taxi Booking Web App & 259 & 47 & 101 & 41 & 57,81\% & Chưa quét & \FAIL \\ \hline
    OFBiz Framework & 3 & 53 & 158 & 97 & 0,96\% & 3 finding & \FAIL \\ \hline
  \end{tabularx}
\end{table}

BookStack đạt Trust Score 98,74\% nhưng vẫn có một component thiếu và một component
thừa nên nội dung là \code{FAIL}. Taxi Booking Web App có đầy đủ các dạng sai lệch,
với Trust Score 57,81\%. OFBiz Framework có chênh lệch lớn giữa source và SBOM đầu
vào; Grype hoàn thành quét và trả về ba finding. Số finding là số cặp CVE--package do
Grype trả về, không nhất thiết bằng số CVE duy nhất và không chứng minh khả năng khai
thác thực tế.

\begin{table}[H]
  \centering
  \caption{Tổng hợp kết quả chức năng}
  \begin{tabularx}{\textwidth}{|C{2.2cm}|Y|C{2.4cm}|L{3.0cm}|}
    \hline
    \rowcolor{headergray}
    \textbf{Mã} & \textbf{Mục tiêu kiểm thử} & \textbf{Luồng chức năng} & \textbf{Nội dung SBOM} \\ \hline
    SBOM-KT01 & Tải lên và phân tích CycloneDX JSON. & \PASS & Không áp dụng \\ \hline
    SBOM-KT02 & Phân tích repository và sinh SBOM bằng Syft. & \PASS & Không áp dụng \\ \hline
    SBOM-KT03 & Đối chiếu SBOM với source và phân loại mismatch. & \PASS & Tùy dữ liệu \\ \hline
    SBOM-KT04 & Tạo pipeline run, cập nhật step và lưu snapshot. & \PASS & Tùy dữ liệu \\ \hline
    SBOM-KT05 & Xuất workbook Excel chứa diff và evidence. & \PASS & Phản ánh run \\ \hline
    SBOM-KT06 & Chạy Grype và phân biệt thành công/thất bại. & \PASS & Enrichment \\ \hline
  \end{tabularx}
\end{table}

\evidencefigure{images/kiem-chung-sbom.png}
  {Màn hình kiểm chứng SBOM và các nhóm sai lệch}

\evidencefigure{images/bao-cao-excel.png}
  {Workbook Excel được xuất từ validation run}

\section{Nhận xét và hạn chế}

\begin{itemize}
  \item Các chức năng cốt lõi của nguyên mẫu đã thực thi được trong phạm vi thử nghiệm.
  \item Trust Score thể hiện độ khớp inventory, không thay thế đánh giá rủi ro bảo mật.
  \item Dữ liệu hiện tại được tạo ở nhiều thời điểm phát triển và chưa được chuẩn hóa
        thành một benchmark cố định.
  \item Grype mới được tích hợp vào các validation run gần đây; các run chưa quét phải
        được hiển thị là ``chưa quét'', không được hiểu là không có CVE.
  \item Dự án chưa có bộ integration test và end-to-end test đầy đủ.
  \item Clone repository, sinh SBOM và quét lỗ hổng đang chạy trong luồng backend;
        repository lớn có thể làm tăng thời gian chờ.
\end{itemize}

\section{Kết luận}

Kết quả thử nghiệm cho thấy hệ thống có thể quản lý, sinh, phân tích và trực quan hóa
SBOM; đồng thời lưu từng validation run, phát hiện các nhóm sai lệch component, tính
Trust Score, bổ sung finding từ Grype và xuất bằng chứng thành Excel. Các kết quả này
xác nhận tính khả thi của nguyên mẫu trong môi trường phát triển, nhưng chưa phải chứng
nhận chất lượng production.

\vspace{1.2cm}
\begin{table}[H]
  \centering
  \begin{tabular}{C{0.45\textwidth}C{0.45\textwidth}}
    \textbf{Người lập tài liệu} & \textbf{Người xác nhận} \\
    \textit{(Ký và ghi rõ họ tên)} & \textit{(Ký và ghi rõ họ tên)} \\
    \\[2.2cm]
    ................................ & ................................
  \end{tabular}
\end{table}

\end{document}
